\section{Single-Call Generation and Safety Verification}
\label{sec:safety}

% Auto-generated by paper/latex/scripts/generate_latex_data.js
\begin{algorithm}[t]
\caption{Interactive Clarification Card Validation}
\label{alg:clarification}
\begin{algorithmic}[1]
\Require User selection $s$, offered candidates $\mathcal{C}$, schema $\mathcal{M}$, bypass flag $b$
\Ensure Validated prompt enrichment string $\mathcal{E}$
\If{$b = \text{true}$}
    \State \Return $\textsc{AppendPromptDirective}(\text{\``User bypassed clarification; do not re-prompt.\'\'}})$
\EndIf
\If{$s = \text{\``Other\'\'}}$
    \State \Return $\textsc{PromptForFreeTextSpecification}()$
\EndIf
\If{$s \notin \mathcal{C} \lor \neg\textsc{ExistsInSchema}(s, \mathcal{M})$}
    \State \Throw \textsc{ValidationError}(\text{HTTP 422, \``Invalid candidate selection\'\'})
\EndIf
\State \Return $\textsc{AppendPromptDirective}(\text{\``Disambiguated target metric: \'\'} + s)$
\end{algorithmic}
\end{algorithm}


\subsection{Card Validation and Prompt Enrichment}
When a user resolves an ambiguity card, Algorithm~\ref{alg:clarification} validates the response before generation. If the user chooses to bypass clarification, the system sets an internal directive preventing downstream re-prompting. If a candidate field is selected, the server verifies that the selection exists both within the offered candidate set and the physical schema allowlist, rejecting tampering with HTTP 422. The disambiguated metric is then appended to the prompt as a structured constraint: ``Disambiguated target metric: $\langle\text{field}\rangle$''.

% Auto-generated by paper/latex/scripts/generate_latex_data.js
\begin{algorithm}[t]
\caption{Recursive Pipeline AST Safety Verification}
\label{alg:safety}
\begin{algorithmic}[1]
\Require Generated pipeline $\mathcal{P} = [\sigma_1, \dots, \sigma_k]$, allowed collection scope $\mathcal{S}$
\Ensure Sanitized pipeline $\mathcal{P}_{\text{safe}}$, classification $\tau \in \{\textsc{ReadOnly}, \textsc{Mutating}\}$
\If{$k > 20$}
    \State \Throw \textsc{SecurityError}(\text{HTTP 422, \``Pipeline length exceeds 20 stages\'\'})
\EndIf
\State $\mathcal{A}_{\text{stages}} \gets \{\texttt{\$match}, \texttt{\$project}, \texttt{\$group}, \texttt{\$sort}, \texttt{\$limit}, \texttt{\$skip}, \texttt{\$unwind}, \texttt{\$lookup}, \dots\}$
\State $\mathcal{D}_{\text{code}} \gets \{\texttt{\$where}, \texttt{\$function}, \texttt{\$accumulator}\}$
\State $\mathcal{W}_{\text{write}} \gets \{\texttt{\$out}, \texttt{\$merge}, \texttt{\$set}, \texttt{\$unset}\}$
\State $\tau \gets \textsc{ReadOnly}$
\For{$\sigma \in \mathcal{P}$}
    \State $k_{\text{stage}} \gets \textsc{StageKey}(\sigma)$
    \If{$k_{\text{stage}} \notin \mathcal{A}_{\text{stages}}$}
        \State \Throw \textsc{SecurityError}(\text{HTTP 422, \``Unallowed stage: \'\'} + $k_{\text{stage}}$)
    \EndIf
    \If{$\textsc{ContainsAnyRecursive}(\sigma, \mathcal{D}_{\text{code}})$}
        \State \Throw \textsc{SecurityError}(\text{HTTP 422, \``Code execution operator forbidden\'\'})
    \EndIf
    \If{$k_{\text{stage}} \in \mathcal{W}_{\text{write}}$}
        \State $\tau \gets \textsc{Mutating}$
    \EndIf
    \State $\mathcal{C}_{\text{ref}} \gets \textsc{ExtractReferencedCollections}(\sigma)$
    \If{$\neg(\mathcal{C}_{\text{ref}} \subseteq \mathcal{S})$}
        \State \Throw \textsc{SecurityError}(\text{HTTP 403, \``Cross-scope collection access forbidden\'\'})
    \EndIf
\EndFor
\If{$\neg\textsc{HasTerminalLimit}(\mathcal{P})$}
    \State $\mathcal{P} \gets \mathcal{P} + [\{\texttt{\$limit}: 1000\}]$
\EndIf
\State \Return $(\mathcal{P}, \tau)$
\end{algorithmic}
\end{algorithm}


\subsection{Recursive Pipeline Safety Guard}
To ensure defense-in-depth, every generated pipeline is parsed into an Abstract Syntax Tree (AST) before database execution (Algorithm~\ref{alg:safety}). The safety validator enforces four strict security invariant layers:
\begin{enumerate}
    \item \textbf{Stage Length Cap:} Pipelines containing more than 20 stages are rejected to mitigate algorithmic resource exhaustion and infinite pipeline denial-of-service.
    \item \textbf{Top-Level and Nested Stage Allowlists:} Only 18 approved analytical aggregation stages (\texttt{\$match}, \texttt{\$project}, \texttt{\$group}, \texttt{\$sort}, \texttt{\$limit}, \texttt{\$skip}, \texttt{\$unwind}, \texttt{\$lookup}, \texttt{\$facet}, \texttt{\$addFields}, \texttt{\$bucket}, \texttt{\$bucketAuto}, \texttt{\$count}, \texttt{\$sortByCount}, \texttt{\$replaceRoot}, \texttt{\$replaceWith}, \texttt{\$sample}, \texttt{\$redact}) are permitted. Furthermore, nested sub-pipelines inside \texttt{\$facet}, \texttt{\$lookup.pipeline}, and \texttt{\$unionWith.pipeline} are recursively subjected to the same allowlist constraints, blocking mutating operations (\texttt{\$out}, \texttt{\$merge}) or administrative stages (\texttt{\$collStats}, \texttt{\$indexStats}, \texttt{\$currentOp}) from being executed inside nested stages.
    \item \textbf{Code Execution Denial:} Arbitrary JavaScript evaluation operators (\texttt{\$where}, \texttt{\$function}, \texttt{\$accumulator}) are categorically blocked at all AST nesting depths.
    \item \textbf{Scope and Document-Form Lookup Verification:} Collections referenced via \texttt{\$lookup}, \texttt{\$unionWith}, or \texttt{\$graphLookup} must reside strictly within the authorized scope $\mathcal{S}$. For document-form lookups (e.g., \texttt{\$lookup: \{ from: \{ coll: '...', db: '...' \} \}}), the database and collection properties are validated separately: referencing administrative namespaces (\texttt{admin}, \texttt{config}, \texttt{local}) or foreign databases outside the connection's target database is blocked to prevent cross-database tenant leaks.
\end{enumerate}

\subsection{Human-in-the-Loop Write Staging}
When a pipeline contains mutating operators (\texttt{\$out}, \texttt{\$merge}, \texttt{\$set}, \texttt{\$unset}), the executor refuses execution and stages the pipeline behind an HMAC-SHA256 signed JSON Web Token (JWT)~\cite{jwt2015rfc}. The token binds the exact pipeline AST, target collection, tenant session, and a unique cryptographic nonce (JTI) valid for ten minutes. Upon client confirmation, the JTI is consumed through a MongoDB unique index constraint, ensuring strict single-use execution and preventing double-spend replay attacks.
